
\subsubsection{2.1 Benchmark Overfitting and Generalization Gaps}

Recht et al. (2019) demonstrated an 11--14 percentage point accuracy gap between ImageNet's original test set and a new independent collection (ImageNet-v2), establishing empirically that models overfit to benchmark-specific features. This cross-sectional approach requires collecting a new independent test set for each benchmark, making it resource-intensive and impossible to apply retrospectively at scale. Engstrom et al. (2020) and Gururangan et al. (2018) examined dataset artifacts as alternative signals, focusing on static structural properties rather than temporal progression. Ethayarajh and Jurafsky (2020) proposed a utility measure for benchmarks but did not address temporal saturation. In contrast, the approach described here requires only historical leaderboard timeseries already available publicly --- no new data collection required.

\subsubsection{2.2 Qualitative Saturation Discourse and Goodhart's Law}

Wang et al. (2019b) explicitly motivated SuperGLUE by GLUE's saturation; Srivastava et al. (2022) similarly motivated BIG-bench by SuperGLUE's saturation. Both treat the saturation event as a community judgment rather than providing a formal detection criterion. Goodhart's (1975) observation --- ''when a measure becomes a target, it ceases to be a good measure'' --- provides the theoretical framing for why benchmarks saturate. Manheim and Garrabrant (2019) operationalize this in the ML context. However, neither work provides a quantitative temporal characterization. Bender et al. (2021) and Kapoor and Narayanan (2023) offer qualitative critiques without automated detection machinery. The gap is a formal, reproducible measurement procedure.

\subsubsection{2.3 Growth Models for Technological Progress}

Logistic growth models have a long history in modeling processes with initial rapid growth followed by a carrying-capacity asymptote: population dynamics \citep{verhulst1838notice}, technology adoption \citep{bass1969new_product}. In the ML context, Kaplan et al. (2020) studied power-law scaling trends but focused on aggregate trends rather than per-benchmark temporal saturation. No prior work is known to apply logistic growth models to benchmark leaderboard timeseries for saturation detection. The AIC framework \citep{burnham2002model} is well-established in ecology; this work transfers it to benchmark dynamics.

\textbf{Positioning}: This work is longitudinal (temporal) rather than cross-sectional; automated and benchmark-agnostic rather than requiring domain-specific new test sets; and formally information-theoretic rather than qualitative.

